% TWC working draft: compile from paper/TWC/.
% Two-column IEEE Transactions journal layout.
\documentclass[10pt,journal]{IEEEtran}
\usepackage[T1]{fontenc}
\usepackage[utf8]{inputenc}
\usepackage{lmodern}
\usepackage{amsmath,amssymb,bm,booktabs,array}
\usepackage{graphicx}
\usepackage{cite}
\usepackage{enumitem}
\usepackage[hidelinks]{hyperref}
\graphicspath{{figures/}}
\allowdisplaybreaks

\newcommand{\CN}{\mathcal{CN}}
\newcommand{\E}{\mathbb{E}}
\newcommand{\I}{\mathbf{I}}
\newcommand{\vect}{\operatorname{vec}}
\newcommand{\tr}{\operatorname{tr}}
\newcommand{\diag}{\operatorname{diag}}
\newcommand{\logit}{\operatorname{logit}}
\newcommand{\sigmoid}{\operatorname{sigmoid}}
\newcommand{\clip}{\operatorname{clip}}
\newcommand{\relu}{\operatorname{Re}}
\newcommand{\imag}{\operatorname{Im}}
\newcommand{\norm}[1]{\left\lVert#1\right\rVert}

\title{Generative-Model-Aided Grant-Free Random Access for GNSS-Independent MIMO-AFDM Satellite IoT Systems}
% TODO: Add the author list and affiliations.
\author{}

\begin{document}
\maketitle
\begin{center}
\small\textit{TWC working draft --- October 21, 2026}
\end{center}
\begin{abstract}
Grant-free random access in low Earth orbit satellite Internet-of-Things
systems must accommodate large differential delays and Doppler shifts
when terminals operate without global navigation satellite system
assistance. These offsets complicate joint activity detection and
channel estimation (JADCE), particularly when fractional delays and
Doppler shifts spread each path's energy across multiple delay--Doppler
grid coefficients. We propose a multi-frame affine frequency division multiplexing
(AFDM) transmission scheme and a deep-unfolded message-passing (MP)
receiver for multiuser multiple-input multiple-output access. We derive
the input--output relationship with integer and fractional offsets and
separate the Doppler-induced phase evolution across frames, allowing
the received pilots to share one effective-channel matrix. We incorporate
a prior learned from multiantenna AFDM channel samples through denoising
score matching into the MP receiver. In each
layer, a channel denoiser and an activity classifier produce the
conditional channel estimate and activity probability, whose product
determines the feedback for subsequent multiuser estimation.
Progressive training and joint fine-tuning optimize the denoisers,
classifiers, and message parameters using received pilots and their
channel and activity labels. This training adapts each layer
to its receiver inputs and to the effects of its estimates on subsequent
interference suppression.
\end{abstract}
% Numerical findings will be added after evaluation of the new receiver.
\begin{IEEEkeywords}
Affine frequency division multiplexing, satellite Internet of Things,
grant-free random access, joint activity detection and channel estimation,
score-based generative models, deep unfolding.
\end{IEEEkeywords}

\section{Introduction}
\label{sec:introduction}

Low Earth orbit (LEO) satellites can complement terrestrial networks
to provide Internet-of-Things (IoT) services in remote areas.
To support access from a large number of IoT devices, grant-free random
access (GFRA) allows active devices to transmit without first obtaining
a grant, thereby reducing signaling overhead and access latency.
However, the receiver has no prior knowledge of which devices are
active. Meanwhile, limited pilot resources make it impractical to assign
mutually orthogonal pilots to all devices. The receiver must therefore
perform joint activity detection and channel estimation (JADCE) from
nonorthogonal pilots \cite{liu2018}. In LEO satellite links, this task
is further complicated by large propagation delays and Doppler shifts.

To compensate for these effects before uplink transmission, the new
radio non-terrestrial network (NR-NTN) framework in Release~17 of the
3rd Generation Partnership Project (3GPP) relies on devices equipped
with a global navigation satellite system (GNSS) receiver. Devices use
their GNSS positions, satellite ephemeris, and the common timing advance
broadcast by the network for delay and Doppler pre-compensation
\cite{3gpp38300}. GNSS hardware and frequent pre-compensation increase
the cost and energy consumption of IoT devices. To support devices
without GNSS receivers or access to GNSS services, Release~19 includes
a study of GNSS-independent operation \cite{3gpp22865}. Handling large
differential delays and Doppler shifts at the satellite receiver reduces
the need for user-specific pre-compensation. The differential delay can
then exceed one symbol duration, while the differential Doppler shift
can exceed one subcarrier spacing. Terrestrial orthogonal frequency
division multiplexing (OFDM) GFRA schemes designed for small timing and
frequency offsets and block-fading channels must therefore be adapted
to these conditions.

Orthogonal time-frequency space (OTFS) modulation supports communication
over doubly selective channels through a delay--Doppler representation
that exploits diversity in both time and frequency \cite{hadani2017}.
For access schemes that handle large offsets at the satellite receiver,
the waveform must also permit simple terminal modulation. Affine
frequency division multiplexing (AFDM) uses a one-dimensional discrete
affine Fourier transform (DAFT) with adjustable chirp parameters.
Properly chosen parameters separate multipath components in the DAFT
domain and enable their diversity to be exploited. AFDM modulation
requires an inverse fast Fourier transform and chirp multiplications,
providing a suitable implementation for terminals with limited
computational resources \cite{bemani2022}.

The channel representation is central to JADCE with AFDM. Conventional
GFRA receivers exploit sparse user activity through compressed sensing
and use channel structure to improve recovery. For example, the
massive multiple-input multiple-output (MIMO) receiver in
\cite{ke2020} jointly exploits activity sparsity and virtual
angular-domain channel sparsity. In GNSS-independent satellite access,
Shen et al.\ \cite{shen2025} combine multi-frame OTFS transmission with
message passing (MP) to jointly recover device activity, channel coefficients, and data
symbols, exploiting channel sparsity in the delay--Doppler--angle domain.
For LEO AFDM GFRA, the graph neural
network in \cite{ding2026} refines a sparse channel prior using nodes
associated with delay--Doppler grid points. Physical delays and Doppler
shifts, however, are continuous. Fractional offsets spread power across
grid coefficients and weaken the assumed sparsity. Increasing the grid
resolution enlarges the dictionary, while estimating continuous offsets
adds path parameters coupled with the unknown gains. Direct estimation
of the AFDM effective channel offers an alternative that incorporates
these offsets into the channel to be recovered.

Score-based generative models learn channel distributions from samples
and use this information for MIMO channel estimation \cite{arvinte2023}.
For AFDM channel estimation, \cite{huang2026} learns the score of the
effective channel and performs posterior inference using the AFDM
observation model. This representation captures fractional-delay and
fractional-Doppler effects without prescribing sparse, grid-aligned
channel support. In multiuser access, the learned channel information
must also support activity detection and interference suppression.
The turbo MP receiver in \cite{cai2026} integrates
score-based channel denoising with linear estimation for MIMO-OFDM
JADCE. It uses the learned prior through local denoising updates,
avoiding online reverse diffusion sampling. These results motivate
combining a learned AFDM effective-channel prior with the joint
estimation structure of MP.

Effective use of this prior also depends on the inputs produced during
message passing. Score pretraining uses Gaussian-perturbed channels,
whereas finite system dimensions and structured sensing matrices can
produce denoiser inputs that depart from the equivalent Gaussian model
\cite{cai2026}. Damping smooths successive messages, but strong damping
can require more iterations to reach a given accuracy
\cite{yoshida2023}. In JADCE, channel and activity errors are coupled
through the feedback: the activity probability weights the channel
estimate used in subsequent interference suppression, and the channel
estimate provides evidence for activity detection. Thus the design
problem includes both adapting the denoiser to its received inputs and
coordinating the two estimators across successive updates.

Deep unfolding turns an iterative receiver into a trainable sequence of
estimation modules \cite{he2020}. We use this structure to optimize
JADCE from received AFDM pilots over a prescribed number of MP layers.
Each layer contains a channel denoiser, an activity classifier, and
trainable message parameters. Joint training propagates channel and
activity losses through their shared feedback: activity supervision
adjusts channel denoising, while channel losses at later layers refine
the activity estimates used by earlier layers. This provides a direct
training mechanism for the interaction between channel recovery,
activity decisions, and multiuser interference suppression.

In this paper, we develop a multi-frame AFDM transmission scheme and a
deep-unfolded MP receiver for GNSS-independent MIMO-AFDM satellite IoT
access. The main contributions are summarized as follows.
\begin{enumerate}[label=\arabic*)]
\item We derive a multi-frame AFDM input--output model for large
differential delays and Doppler shifts, retaining both integer and
fractional offsets. Separating the inter-frame Doppler phase evolution
places these phase factors in the pilot operator and allows all pilot
frames within an access block to share an effective-channel matrix.
The resulting formulation combines additional pilot observations
without introducing a separate set of channel coefficients for each
frame.

\item We develop a channel and activity estimation module based on a
learned multiantenna AFDM effective-channel prior. The score-based
denoiser estimates each user's channel conditional on being active, and a
trainable classifier estimates its activity probability from the
denoiser input and channel estimate. Their product is used to form the
feedback to the linear estimator, integrating off-grid channel
structure and soft activity information into multiuser MP.

\item We unfold the MP receiver into layers with independently
trainable denoisers, activity classifiers, and message parameters.
A progressive training procedure adds layers and jointly fine-tunes
the receiver using channel and activity supervision. By propagating
both losses through the inter-layer messages, the training adapts the
two estimators to the evolving receiver inputs and their effects on
subsequent interference suppression.
\end{enumerate}

% Numerical findings will be added after independent evaluation of the new receiver.
Section~\ref{sec:system} presents the system model.
Sections~\ref{sec:algorithm}--\ref{sec:overall} develop the receiver,
its training procedure, and computational complexity. Sections~VI and
VII present numerical results and conclusions, respectively.


\noindent\textit{Notation:} $\CN(\bm\mu,\bm\Sigma)$ denotes a proper complex Gaussian law.
$\|\cdot\|_F$ denotes the Euclidean norm of all matrix or tensor entries;
$(\cdot)^{\mathsf H}$ and $\tr(\cdot)$ denote conjugate transpose and trace.

\section{System Model}
\label{sec:system}

We consider a grant-free uplink from $K$ single-antenna terminals to a
low Earth orbit (LEO) satellite equipped with a uniform planar array (UPA).
The array has $M=M_yM_z$ elements. An access block contains $Q$ affine
frequency division multiplexing (AFDM) pilot frames, each with $N$ useful
samples. Terminal activities, channel geometry, and path coefficients
remain fixed within the block.
The receiver knows the pilot sequences and the system parameters,
whereas terminal activities, channel realizations, inter-frame phase
increments, and receiver noise power are unknown.

Let $\alpha_k\in\{0,1\}$ denote the activity state of terminal $k$.
In the activity model, independent activations with
probability $p$ are conditioned on at least one transmission:
\begin{equation}
 \begin{aligned}
 p_{\rm gen}(\bm\alpha)&=
 \frac{\mathbf{1}\{\sum_{k=1}^{K}\alpha_k\geq1\}}{Z}
 \prod_{k=1}^{K}p^{\alpha_k}(1-p)^{1-\alpha_k},\\
 Z&=1-(1-p)^K.
 \end{aligned}
 \label{eq:activity-generation}
\end{equation}
The resulting activity marginal is $\rho=p/Z$, which is supplied to
the receiver.

\subsection{Terrestrial-Satellite Link Model}

Let $\mathbf{s}$ and $\mathbf{v}_{\rm sat}$ denote the satellite position
and velocity, respectively. The position of terminal $k$ is denoted by
$\mathbf{g}_k$.
With carrier frequency $f_c$ and propagation speed $c_0$, define
\begin{equation}
 \mathbf{u}_k=\frac{\mathbf{g}_k-\mathbf{s}}
 {\norm{\mathbf{g}_k-\mathbf{s}}},
 \qquad
 \nu_k=-\frac{f_c}{c_0}\mathbf{v}_{\rm sat}^{\mathsf T}\mathbf{u}_k.
 \label{eq:geometry}
\end{equation}
We use the complex-envelope convention
$s_{\rm RF}(t)=\Re\{x(t)e^{-j2\pi f_ct}\}$; thus, a path contributes
$x(t-\tau(t))e^{j2\pi f_c\tau(t)}$ up to its complex gain.
Equation~\eqref{eq:geometry} consequently gives the Doppler time factor
$e^{j2\pi\nu_kt}$. Envelope delays and array directions are held fixed
over the block, while this carrier phase variation is retained.

After compensation of the common reference delay $\tau_c$, the residual
delay of path $\ell$ is given by
\begin{equation}
 \tau_{k\ell}=\frac{\norm{\mathbf{g}_k-\mathbf{s}}}{c_0}-\tau_c
             +\delta_\ell\tau_{{\rm spread},k},
 \label{eq:path-delay}
\end{equation}
where $\tau_c$ is the propagation time to the closest ground point in the
beam footprint, and $\delta_\ell\tau_{{\rm spread},k}$ is the excess
multipath delay. We adopt a common spatial response and a common Doppler
shift for the paths of each terminal
\cite{li2023}\cite{huang2026}.

Let $\mu_{y,k}=\mathbf e_y^{\mathsf T}\mathbf u_k$ and
$\mu_{z,k}=\mathbf e_z^{\mathsf T}\mathbf u_k$, where $\mathbf e_y$ and
$\mathbf e_z$ are the array axes. For element position
$\mathbf s-n_yd_y\mathbf e_y-n_zd_z\mathbf e_z$, the array response is
\begin{equation}
 \begin{gathered}
 [\mathbf{a}_k]_m=
 \exp\!\left[j2\pi\left(
 \frac{d_y}{\lambda}n_y\mu_{y,k}
 +\frac{d_z}{\lambda}n_z\mu_{z,k}\right)\right],\\
 m=n_z+M_zn_y,\qquad \norm{\mathbf{a}_k}^2=M,
 \end{gathered}
 \label{eq:upa}
\end{equation}
where $\lambda=c_0/f_c$, $0\leq n_y<M_y$, and $0\leq n_z<M_z$.
The unknown array response is included in the channel to be estimated.

\subsection{Transmission Scheme}

Let $T_s=1/(N\Delta f)$ be the sampling interval, where $\Delta f$
denotes the subcarrier spacing. The discrete affine Fourier transform
(DAFT) matrix is \cite{bemani2022}
\begin{equation}
 \begin{gathered}
 [\mathbf{F}_{\rm A}]_{rt}=\frac{1}{\sqrt N}
 e^{-j2\pi(c_2r^2+rt/N+c_1t^2)},\\
 r,t=0,\ldots,N-1.
 \end{gathered}
 \label{eq:daft}
\end{equation}
For pilot vector $\mathbf p_{qk}$, the useful time samples are
$\mathbf s_{qk}=\mathbf F_{\rm A}^{\mathsf H}\mathbf p_{qk}$.
A chirp-periodic prefix (CPP) of length $N_{\rm CPP}$ extends these samples as
\begin{equation}
 \begin{gathered}
 s_{qk}[t]=e^{j2\pi c_1(t^2-\langle t\rangle_N^2)}
 s_{qk}[\langle t\rangle_N],\\
 -N_{\rm CPP}\leq t<0,
 \end{gathered}
 \label{eq:cpp}
\end{equation}
where $\langle x\rangle_N=x-N\lfloor x/N\rfloor$.
We choose $N_{\rm CPP}T_s\geq\max_{k,\ell}\tau_{k\ell}$ and
$2Nc_1\in\mathbb N_+$. The frame duration is
$T_{\rm frame}=(N+N_{\rm CPP})T_s$.

To account for continuous delays, we use the AFDM waveform model in
\cite{huang2026}. Let $w_\xi=\langle\xi\rangle_N$ and
$\zeta_n(w)=\lfloor2c_1w+n/N\rfloor$. The basis is
\begin{equation}
 \varphi_n(\xi)=\frac{1}{\sqrt N}
 e^{j2\pi[c_2n^2+c_1\xi^2+nw_\xi/N-\zeta_n(w_\xi)w_\xi]}.
 \label{eq:continuous-waveform}
\end{equation}
It agrees with the inverse DAFT at integer useful-sample times and
extends across prefix periods according to \eqref{eq:cpp}.
For a delay $d=\tau/T_s$, set $w=\langle t-d\rangle_N$ and
$\zeta_{tn}=\zeta_n(w)$. The delayed basis can be expressed as
\begin{equation}
 \begin{aligned}
 [\mathbf B_d]_{tn}&=\varphi_n(t-d)\\
 &=\frac{1}{\sqrt N}
 e^{j2\pi[c_2n^2+c_1(t-d)^2+nw/N-\zeta_{tn}w]}.
 \end{aligned}
 \label{eq:continuous-basis}
\end{equation}
The corresponding DAFT-domain path operator, for useful-window start
index $g$, is given by
\begin{equation}
 \begin{aligned}
 \mathbf G(d,\nu;g)&=\mathbf F_{\rm A}\mathbf D_{\nu,g}\mathbf B_d,\\
 \mathbf D_{\nu,g}&=\diag\{e^{j2\pi\nu T_s(g+t)}\}_{t=0}^{N-1}.
 \end{aligned}
 \label{eq:path-operator}
\end{equation}
Expanding the path operator gives
\begin{equation}
 \begin{split}
 G_{rn}
 &=\frac{e^{j2\pi[c_1d^2+c_2(n^2-r^2)-nd/N+\nu T_sg]}}{N}\\
 &\quad\times\sum_{t=0}^{N-1}
 e^{j2\pi[(n-r)/N-2c_1d+\nu T_s]t}
 e^{j2\pi\zeta_{tn}d}.
 \end{split}
 \label{eq:path-entry-expansion}
\end{equation}
The last factor accounts for the fractional delay in the waveform.
This expression corresponds to \cite{huang2026} with normalized
Doppler $k_{\nu,\mathrm{Huang}}=-\nu NT_s$ under our sign convention.

For integer $d$, \eqref{eq:path-operator} reduces to
\begin{align}
 \mathbf G(d,\nu;g)
 &=\mathbf F_{\rm A}\mathbf D_{\nu,g}\bm\Omega_d
   \bm\Pi_{d\bmod N}\mathbf F_{\rm A}^{\mathsf H},
 \label{eq:integer-operator}\\
 [\bm\Omega_d]_{tt}
 &=e^{j2\pi c_1[(t-d)^2-\langle t-d\rangle_N^2]},
 \label{eq:integer-prefix-phase}
\end{align}
where $[\bm\Pi_l\mathbf s]_t=s[\langle t-l\rangle_N]$.
The full delay enters the prefix phase, and its remainder determines
the cyclic shift.

At the reference frame, the channel of terminal $k$ is given by
\begin{equation}
 \begin{aligned}
 \mathbf G_k&=\sum_\ell g_{k\ell}
 \mathbf G(\tau_{k\ell}/T_s,\nu_k;N_{\rm CPP}),\\
 H_{kmrn}&=[\mathbf a_k]_m[\mathbf G_k]_{rn},
 \end{aligned}
 \label{eq:channel}
\end{equation}
where $g_{k\ell}$ includes the path gain after large-scale power
compensation. The tensor $H_k\in\mathbb C^{M\times N\times N}$ contains
all antenna, DAFT output, and DAFT input coordinates. Reshaping it as a matrix
in $\mathbb C^{M\times N^2}$ with one row per antenna yields a matrix of rank at most one under
the common spatial response assumption.

Pilots are formed from selected rows and permuted columns of a
$KN$-point discrete Fourier transform (DFT). With distinct row indices
$r_q$, a column permutation $c_{kn}$, and independent uniform phases
$\phi_{kn}$, set
\begin{equation}
 \begin{gathered}
 P_{qkn}=\sqrt{\frac{E_p}{N}}
 e^{-j2\pi r_qc_{kn}/(KN)}e^{j\phi_{kn}},\\
 \sum_{n=0}^{N-1}|P_{qkn}|^2=E_p.
 \end{gathered}
 \label{eq:pilot}
\end{equation}
The common terminal Doppler introduces the phase increment
$\beta_k=\operatorname{wrap}(2\pi\nu_kT_{\rm frame})$ between consecutive
frames, where $\operatorname{wrap}$ maps an angle to $[-\pi,\pi)$.
The observation operator with these phase increments is
\begin{equation}
 A_{qkn}=P_{qkn}e^{jq\beta_k},\qquad q=0,\ldots,Q-1.
 \label{eq:effective-pilot}
\end{equation}

Define the effective channel $X_k=\alpha_kH_k$.
After prefix removal and DAFT demodulation, the received signal can be expressed as
\begin{equation}
 \begin{gathered}
 Y_{qmr}=\sum_{k=1}^K\sum_{n=0}^{N-1}A_{qkn}X_{kmrn}+W_{qmr},
 \\
 W_{qmr}\sim\CN(0,\sigma^2),
 \end{gathered}
 \label{eq:observation-tensor}
\end{equation}
with conditionally independent noise entries.
Using the indexing
\begin{equation}
 \begin{aligned}
 [\mathbf Y]_{q,mN+r}&=Y_{qmr},\\
 [\mathbf A]_{q,(k-1)N+n}&=A_{qkn},\\
 [\mathbf X]_{(k-1)N+n,mN+r}&=X_{kmrn},
 \end{aligned}
 \label{eq:mmv-indexing}
\end{equation}
the multiple measurement vector (MMV) model can be written as
\begin{equation}
 \begin{gathered}
 \mathbf Y=\mathbf A\mathbf X+\mathbf W,\\
 \mathbf Y\in\mathbb C^{Q\times MN},\qquad
 \mathbf A\in\mathbb C^{Q\times KN},\\
 \mathbf X\in\mathbb C^{KN\times MN}.
 \end{gathered}
 \label{eq:mmv}
\end{equation}
All $MN$ measurement vectors share the terminal activity pattern.
The complete receiver takes the observation tensor, the original pilots $P$,
and the activity marginal $\rho$ as inputs. It estimates
$\widehat{\bm\beta}$ and $\widehat\sigma^2$ from these observations
and constructs $\widehat{\mathbf A}=\mathbf A(\widehat{\bm\beta})$;
the estimation procedure is given in Appendix~\ref{app:frontend}.
The physical operator $\mathbf A=\mathbf A(\bm\beta)$ and the noise
variance $\sigma^2$ in \eqref{eq:mmv} describe data generation.
For simulations, the signal-to-noise ratio (SNR)
is defined using the clean superposition $\mathbf Y_0=\mathbf A\mathbf X$:
\begin{equation}
 \sigma^2=\frac{\norm{\mathbf Y_0}_F^2}{QMN}
 10^{-\mathrm{SNR}_{\rm dB}/10}.
 \label{eq:snr}
\end{equation}
The received signal power includes cross terms between terminals.

Write $(\mathbf A_{\rm c},\sigma_{\rm c}^2)$ for the operator and noise
variance held fixed during the message updates. For the
complete receiver these are $(\widehat{\mathbf A},\widehat\sigma^2)$;
in controlled training they are the supplied $(\mathbf A,\sigma^2)$.

\subsection{Problem Formulation}

Let $D=MN^2$ and let $q_{\rm tr}$ be the distribution of normalized channels
generated from the physical model before applying activity masks.
We normalize channels and observations by a common scale $s$,
the root mean square of the training channels:
$(\mathbf Y,H_k,X_k,\mathbf W)$ are divided by $s$, and both
$\sigma^2$ and $\sigma_{\rm c}^2$ by $s^2$. The pilot operators are
unchanged. The same symbols denote the normalized quantities below;
the reported channel estimates are rescaled by $s$.

For joint activity detection and channel estimation (JADCE), we use
the factorized prior
\begin{equation}
 q(\{dH_k\},\bm\alpha)=
 \prod_{k=1}^{K}q_{\rm tr}(dH_k)
 \rho^{\alpha_k}(1-\rho)^{1-\alpha_k}.
 \label{eq:working-prior}
\end{equation}
We approximate the joint prior by the product of its channel
and activity marginals. With the sensing matrix and noise variance
held fixed, the approximate posterior distribution can be expressed as
\begin{equation}
 \begin{aligned}
 &q(d\mathbf X,\{dH_k\},\bm\alpha\mid
 \mathbf Y,\mathbf A_{\rm c},\sigma_{\rm c}^2)\\
 &\quad\propto
 e^{-\norm{\mathbf Y-\mathbf A_{\rm c}\mathbf X}_F^2/\sigma_{\rm c}^2}\\
 &\qquad\times\prod_{k=1}^{K}
 \delta_{\alpha_kH_k}(dX_k)q_{\rm tr}(dH_k)
 \rho^{\alpha_k}(1-\rho)^{1-\alpha_k},
 \end{aligned}
 \label{eq:working-posterior}
\end{equation}
where $\delta_{\alpha_kH_k}$ enforces $X_k=\alpha_kH_k$.
Under this posterior, the minimum mean-squared error (MMSE) target
and the activity decision are given by
\begin{equation}
 \begin{aligned}
 X_k^{\rm MMSE}&=\E_q[X_k\mid\mathbf Y,\mathbf A_{\rm c},\sigma_{\rm c}^2],\\
 p_k^{\rm Bayes}&=q(\alpha_k=1\mid\mathbf Y,\mathbf A_{\rm c},\sigma_{\rm c}^2),\\
 \alpha_k^{\rm Bayes}&=\mathbf1\{p_k^{\rm Bayes}\geq\tfrac12\}.
 \end{aligned}
 \label{eq:targets}
\end{equation}
The threshold corresponds to equal costs for a missed detection and a
false alarm. The receiver approximates these targets by iterating between
linear estimation and channel and activity denoising, returning
$\widehat X_k$, $p_k$, and $\widehat\alpha_k=\mathbf1\{p_k\geq1/2\}$.

\section{Deep-Unfolded Receiver Design}
\label{sec:algorithm}

We develop a deep-unfolded MP receiver with a score-based channel
prior for joint activity detection and channel estimation. The
likelihood and priors in \eqref{eq:working-posterior} provide a
starting point for constructing its estimation modules. Following
\cite{cai2026}, a linear estimator processes the received pilots,
while a channel and activity denoiser incorporates the prior
information; the two modules exchange Gaussian extrinsic messages. To adapt the denoisers to the inputs
produced by this receiver, we unfold the message updates into $L$
layers with trainable channel denoisers, activity classifiers, and
message parameters. The linear estimator uses the sensing matrix
formed from the pilots and inter-frame phase increments. The nonlinear
module estimates each user's conditional channel and activity
probability, whose product determines the feedback to the next layer.
Section~\ref{sec:learning} trains these coupled updates using received
pilot observations. In the following,
$J=KN$ and $C=MN$.

\subsection{Linear Estimation}

At layer $\ell$, the input Gaussian messages have means collected in
$\mathbf M_\ell$ and a common covariance $w_\ell\I_J$ for each
column of $\mathbf X$. Combining them with the likelihood in
\eqref{eq:working-posterior} and completing the square gives the
mean $\mathbf M_\ell+\bm\Delta_\ell$ and average marginal variance
$t_\ell$, where
\begin{equation}
 \begin{aligned}
 \bm\Delta_\ell&=w_\ell\mathbf A_{\rm c}^{\mathsf H}
 (\sigma_{\rm c}^2\I_Q+w_\ell\mathbf A_{\rm c}\mathbf A_{\rm c}^{\mathsf H})^{-1}
 (\mathbf Y-\mathbf A_{\rm c}\mathbf M_\ell),
 \\
 t_\ell&=\frac1J\tr\!\left[
 \left(w_\ell^{-1}\I_J+
 \frac{\mathbf A_{\rm c}^{\mathsf H}\mathbf A_{\rm c}}{\sigma_{\rm c}^2}
 \right)^{-1}\right].
 \end{aligned}
 \label{eq:twc-linear-correction}
\end{equation}
Approximating the covariance by $t_\ell\I_J$, we obtain the
extrinsic message by dividing the resulting Gaussian by the input
message \cite{minka2001}. Its variance and mean are given by
\cite{cai2026}
\begin{equation}
 \begin{aligned}
 v_\ell^{(0)}&=(t_\ell^{-1}-w_\ell^{-1})^{-1},\\
 \mathbf R_\ell^{(0)}&=v_\ell^{(0)}
 \left(\frac{\mathbf M_\ell+\bm\Delta_\ell}{t_\ell}
       -\frac{\mathbf M_\ell}{w_\ell}\right).
 \end{aligned}
 \label{eq:twc-linear-extrinsic}
\end{equation}
To adapt these messages during receiver training, we introduce a
positive factor $\gamma_\ell=e^{c_{\gamma,\ell}}$ that scales the
extrinsic variance and adjust the mean as
\begin{equation}
 v_\ell=\gamma_\ell v_\ell^{(0)},\qquad
 \mathbf R_\ell=\mathbf M_\ell+
 \left(1+\frac{v_\ell}{w_\ell}\right)\bm\Delta_\ell.
 \label{eq:twc-query-calibrated}
\end{equation}
For $\gamma_\ell=1$, \eqref{eq:twc-query-calibrated} reduces to
\eqref{eq:twc-linear-extrinsic}. The pair
$(\mathbf R_\ell,v_\ell)$ is then passed to the channel and activity
denoiser.

\subsection{Channel and Activity Denoising}

The Gaussian messages define a local estimation model for constructing
the channel denoiser and activity classifier. For a user block
$r\in\mathbb C^D$ of $\mathbf R_\ell$, the input message is
proportional to $\CN(r;\alpha h,v\I_D)$.
Combining this message with the channel prior in
\eqref{eq:working-prior}, the densities of $r$ under the inactive
and active hypotheses are, respectively,
\begin{equation}
 \begin{aligned}
 q_0(r;v)&=\CN(r;0,v\I_D),\\
 q_1(r;v)&=\int\CN(r;h,v\I_D)\,q_{\rm tr}(dh).
 \end{aligned}
 \label{eq:twc-local-densities}
\end{equation}
Under this local model, the joint channel and activity posterior is
\begin{equation}
 \begin{gathered}
 q(dh,a\mid r,v)=
 \frac{\rho^a(1-\rho)^{1-a}\CN(r;ah,v\I_D)\,q_{\rm tr}(dh)}
 {(1-\rho)q_0(r;v)+\rho q_1(r;v)},\\
 a\in\{0,1\}.
 \end{gathered}
 \label{eq:twc-local-posterior}
\end{equation}
Conditioning \eqref{eq:twc-local-posterior} on the active hypothesis
gives the conditional channel mean
\begin{equation}
 \mu^\star(r,v)=\E[h\mid r,\alpha=1]
 =\frac{\int h\,\CN(r;h,v\I_D)\,q_{\rm tr}(dh)}{q_1(r;v)}.
 \label{eq:twc-local-channel-target}
\end{equation}
The activity probability follows from Bayes' rule as
\begin{equation}
 p^\star(r,v)=\Pr(\alpha=1\mid r)
 =\frac{\rho q_1(r;v)}
 {(1-\rho)q_0(r;v)+\rho q_1(r;v)}.
 \label{eq:twc-local-activity-target}
\end{equation}
Since $x=\alpha h$ vanishes under the inactive hypothesis, averaging
over the two activity states gives \cite{cai2026}
\begin{equation}
 \E[x\mid r]=p^\star(r,v)\mu^\star(r,v).
 \label{eq:twc-mixture-mean}
\end{equation}
Thus channel denoising and activity detection provide the two
quantities needed to estimate the effective channel. Computing them
directly requires integration over the $D$-dimensional complex channel
distribution in \eqref{eq:twc-local-densities} and
\eqref{eq:twc-local-channel-target}. This decomposition motivates a
score-based channel denoiser and an activity classifier whose outputs
are multiplied to form the effective-channel estimate. Their
parameters are then jointly trained through the unfolded iterations
using channel and activity supervision.

Let $u_{k\ell}\in\mathbb R^{2D}$ denote the real and imaginary parts
of the denoiser input for user $k$. Based on the score representation developed
in Section~\ref{sec:learning}, the channel estimate is
\begin{equation}
 \mu_{k\ell}=u_{k\ell}
       +a_\ell\frac{v_\ell}{2}s_{\theta_\ell}(u_{k\ell},v_\ell),
 \qquad a_\ell=e^{c_{\alpha,\ell}}.
 \label{eq:twc-conditional-mean}
\end{equation}
Here $s_{\theta_\ell}$ is the score network in layer $\ell$,
$v_\ell/2$ is the noise variance per real coordinate, and $a_\ell$
controls the denoising correction.

For the classifier, rewriting \eqref{eq:twc-local-activity-target}
in terms of posterior odds gives
\begin{equation}
 \logit p^\star(r,v)
 =\logit(\rho)+\log\frac{q_1(r;v)}{q_0(r;v)}.
 \label{eq:twc-activity-odds}
\end{equation}
The log density ratio supplies the observation-dependent correction
to the prior log odds under the local model. Motivated by this form,
we parameterize the correction with a trainable classifier
\cite{menon2016}.
Using the denoiser input and channel estimate, we compute
\begin{equation}
 p_{k\ell}=\sigmoid\!\left(\logit(\rho)+
 g_{\phi_\ell}(u_{k\ell},v_\ell,\mu_{k\ell},\rho)\right),
 \label{eq:twc-activity}
\end{equation}
where $g_{\phi_\ell}$ parameterizes this correction. Substituting
the channel and activity estimates into
\eqref{eq:twc-mixture-mean} gives
\begin{equation}
 \widehat x_{k\ell}=p_{k\ell}\mu_{k\ell}.
 \label{eq:twc-estimate}
\end{equation}
The user estimates are converted back to complex form and arranged
as $\widehat{\mathbf X}_\ell$ according to \eqref{eq:mmv-indexing}.
Within each layer, all users share the score network and activity
classifier.

\subsection{Extrinsic Message Update}

To form the message returned to the linear estimator, we combine
the mean $\widehat{\mathbf X}_\ell$ with a shared variance estimated
from the observation residual. This supplies the feedback variance
without a second-order score network. We subtract the noise
contribution from the squared residual and normalize it as
\begin{equation}
 e_\ell=\frac{\|\mathbf Y-\mathbf A_{\rm c}\widehat{\mathbf X}_\ell\|_F^2
                   -QC\sigma_{\rm c}^2}
                  {C\|\mathbf A_{\rm c}\|_F^2}.
 \label{eq:twc-residual-scale}
\end{equation}
We then compute the ratio $q_\ell$ of the output variance to the input
variance through
\begin{equation}
 \begin{aligned}
 z_\ell&=e^{c_{\beta,\ell}}
 \left[\eta v_\ell\operatorname{softplus}
       \!\left(\frac{e_\ell}{\eta v_\ell}\right)+\epsilon_s\right],\\
 q_\ell&=\epsilon_r+(q_{\max}-\epsilon_r)
 \tanh\!\left(\frac{z_\ell}{(q_{\max}-\epsilon_r)v_\ell}\right),
 \end{aligned}
 \label{eq:twc-feedback-ratio}
\end{equation}
where $\operatorname{softplus}(x)=\log(1+e^x)$,
$\eta,\epsilon_s>0$, and $0<\epsilon_r<q_{\max}<1$ are fixed.
Subtracting the noise contribution can make $e_\ell$ negative;
the softplus transformation makes it positive and differentiable
for training.
The learned multiplier $e^{c_{\beta,\ell}}$ scales the positive
quantity in brackets, while the mapping bounds the output variance
$v_\ell q_\ell$ below the input variance $v_\ell$. This upper bound
ensures a positive variance after Gaussian message division.
Dividing this Gaussian output message by the input message gives
\begin{equation}
 \widetilde{\mathbf M}_{\ell+1}
 =\frac{\widehat{\mathbf X}_\ell-q_\ell\mathbf R_\ell}{1-q_\ell},
 \qquad
 \widetilde w_{\ell+1}=\frac{v_\ell q_\ell}{1-q_\ell}.
 \label{eq:twc-feedback-proposal}
\end{equation}
We damp the resulting message in its natural parameters
\cite{vehtari2020}:
\begin{equation}
 \begin{aligned}
 w_{\ell+1}^{-1}&=(1-\delta_\ell)w_\ell^{-1}
                    +\delta_\ell\widetilde w_{\ell+1}^{-1},\\
 \frac{\mathbf M_{\ell+1}}{w_{\ell+1}}
 &=(1-\delta_\ell)\frac{\mathbf M_\ell}{w_\ell}
       +\delta_\ell\frac{\widetilde{\mathbf M}_{\ell+1}}
                           {\widetilde w_{\ell+1}}.
 \end{aligned}
 \label{eq:twc-damping}
\end{equation}
Here $\delta_\ell=\min\{1,w_{\rm ref}/\widetilde w_{\ell+1}\}$
and $w_{\rm ref}=\rho P_H$, with
$P_H=\E_{q_{\rm tr}}\|h\|^2/D$ denoting the average channel power
per coefficient.
The damping reduces the weight of an outgoing message when its
variance exceeds $w_{\rm ref}$. Equations~\eqref{eq:twc-feedback-proposal}
and \eqref{eq:twc-damping} complete the update to layer $\ell+1$.
Through this feedback, both channel and activity estimates affect
the inputs to subsequent layers. The next section develops the
training objective for these coupled updates.

\section{Score Pretraining and Joint Receiver Training}
\label{sec:learning}

Training the deep-unfolded MP receiver in Section~\ref{sec:algorithm}
starts with learning a score-based channel prior. Score pretraining
uses Gaussian-perturbed samples of multiantenna AFDM effective
channels generated with continuous path parameters. The pretrained
network initializes the channel denoisers. We then jointly optimize
the denoisers, activity classifiers, and message parameters using
received pilot observations, allowing channel and activity errors
to guide the updates across receiver layers.

\subsection{First-Order Score Matching}

Let $h_R\in\mathbb R^{2D}$ denote a normalized channel sample.
We perturb it as $u=h_R+\sqrt{v/2}\,\varepsilon$, where
$\varepsilon\sim\mathcal N(0,\I_{2D})$. Denote the perturbed
density by $q_{1,R}(u;v)$ and its score by
$s^\star(u,v)=\nabla_u\log q_{1,R}(u;v)$.
Tweedie's formula relates this score to the conditional channel mean
\cite{efron2011,cai2026}:
\begin{equation}
 \E[h_R\mid u]=u+\frac v2\,s^\star(u,v).
 \label{eq:twc-tweedie}
\end{equation}
Equation~\eqref{eq:twc-tweedie} expresses the conditional mean in
\eqref{eq:twc-local-channel-target} through the score of the perturbed
channel distribution. Replacing this score by a network gives the initial denoising
correction in \eqref{eq:twc-conditional-mean} with $a_\ell=1$.
Joint receiver training subsequently optimizes the network and this
coefficient for channel estimation through the unfolded layers.

To learn the score over the input noise levels, consider the
weighted score matching objective \cite{cai2026}
\begin{equation}
 \mathcal L_{\rm SM}(\theta)
 =\E_{v,u}\!\left[
 \frac{v}{4D}\left\|s_\theta(u,v)
       -\nabla_u\log q_{1,R}(u;v)\right\|^2\right],
 \label{eq:twc-score-matching}
\end{equation}
where $u$ follows $q_{1,R}(\cdot;v)$ at each noise level.
The marginal score in this objective is unknown. Denoising score
matching replaces it by the conditional perturbation score and has
the same minimizers \cite{vincent2011,song2019}.
The Gaussian perturbation has conditional score
$\nabla_u\log q(u\mid h_R,v)=-\varepsilon/\sqrt{v/2}$.
Substitution gives the training loss
\begin{equation}
 \mathcal L_{\rm DSM}(\theta)=\E_{h_R,v,\varepsilon}
 \left[\frac{\|\sqrt{v/2}\,s_\theta(h_R+\sqrt{v/2}\varepsilon,v)
                      +\varepsilon\|^2}{2D}\right].
 \label{eq:twc-dsm}
\end{equation}
Training samples are drawn from the channel distribution before
activity masking and perturbed over a range of noise levels.
The resulting network initializes $s_{\theta_\ell}$ in
\eqref{eq:twc-conditional-mean}. During joint receiver training, this
score-based form is retained while its network weights and denoising
coefficient are optimized for the inputs produced by the unfolded
layers.

\subsection{Network Inputs}

For each terminal, the real and imaginary antenna maps are arranged
as $2M$ channels over the $N\times N$ DAFT coordinates. The score
network is conditioned on $v_\ell$ and returns a tensor of the same
shape, which is flattened to the vector used in
\eqref{eq:twc-conditional-mean}. This representation retains the full
multiantenna channel within each network evaluation.

The activity classifier concatenates the denoiser input and its
channel estimate into $4M$ input channels and scales them by
$(1+v_\ell)^{-1/2}$. Three convolutional layers followed by spatial
average pooling extract a feature vector. Two fully connected layers
combine this vector with the five scalar inputs
\begin{equation}
 \begin{aligned}
 \Big(&\log v_\ell,\
 \log\!\left(1+\frac{\|u_{k\ell}\|^2}{Dv_\ell}\right),\
 \log\!\left(1+\frac{\|\mu_{k\ell}\|^2}{Dv_\ell}\right),\\
 &\log\!\left(1+\frac{\|u_{k\ell}-\mu_{k\ell}\|^2}{Dv_\ell}\right),\
 \logit\rho\Big)
 \end{aligned}
 \label{eq:twc-classifier-features}
\end{equation}
to produce the log-odds correction in \eqref{eq:twc-activity}.
The energy terms describe the denoiser input, channel estimate, and denoising
residual relative to the input noise level. Each network acts on
one terminal at a time, with parameters shared across terminals
within a layer. Thus its input dimensions depend on $M$ and $N$,
while $K$ determines the number of evaluations and $Q$ enters through
the linear estimator.

\subsection{Joint Channel and Activity Training}

To adapt the denoisers to the inputs produced by the receiver, we
pass each training observation through the updates in
Section~\ref{sec:algorithm}. In this training stage, the sensing matrix
and noise variance are supplied from the physical model, as specified
in Section~\ref{sec:system}; phase and noise estimation are performed
separately in the complete receiver.
The denoiser inputs and feedback messages depend on the current layer
parameters. We supervise the two factors in
\eqref{eq:twc-mixture-mean} using the squared channel error on active
users and binary cross-entropy on all users. These losses provide
distinct supervision for the conditional channel estimate and activity
probability within a common objective. Their
product determines the feedback to later layers, so both estimators
participate in joint optimization.

Suppose the receiver training set contains $B$ scenes, with $K_i$
candidate users in scene $i$. Let $h_{ik,R}$ and $\alpha_{ik}$
denote the channel and activity label of user $k$ in scene $i$.
Training and validation use disjoint physical scenes.
The channel and activity losses at layer $\ell$ are respectively
given by
\begin{align}
 R_H^{(\ell)}&=\frac{1}{E_{\rm tr}}\sum_{i,k}
       \alpha_{ik}\|\mu_{ik\ell}-h_{ik,R}\|^2,
 \label{eq:twc-channel-loss}\\
 R_a^{(\ell)}&=-\frac{1}{T_{\rm tr}}\sum_{i,k}
       \left[\alpha_{ik}\log p_{ik\ell}
       +(1-\alpha_{ik})\log(1-p_{ik\ell})\right].
 \label{eq:twc-activity-loss}
\end{align}
Here $E_{\rm tr}=\sum_{i,k}\alpha_{ik}\|h_{ik,R}\|^2$ and
$T_{\rm tr}=\sum_i K_i$ are the fixed active-channel energy and user
count of the training set. The outputs $\mu_{ik\ell}$ and $p_{ik\ell}$
are computed from the received observations; the activity labels
select the channel terms in \eqref{eq:twc-channel-loss}.

For a receiver with $m$ layers, the joint training objective is
\begin{equation}
 J_m=\frac{\sum_{\ell=1}^{m}\omega_\ell
                  \left(R_H^{(\ell)}+R_a^{(\ell)}\right)}
                 {\sum_{\ell=1}^{m}\omega_\ell},
 \label{eq:twc-layer-objective}
\end{equation}
where $\omega_\ell$ is the weight assigned to layer $\ell$.
The last layer provides the final estimates, while earlier-layer
outputs provide intermediate supervision.
Gradients of \eqref{eq:twc-layer-objective} are propagated through
the denoisers, classifiers, and inter-layer messages. Thus activity
supervision also adjusts channel denoising, and channel losses at later
layers refine the activity estimates used in earlier feedback.

\subsection{Progressive Training}

We increase receiver depth progressively. Each new layer has independent
parameters, with its score network initialized from the same
pretrained network.
Its classifier output $g_{\phi_\ell}$ and the parameters $c_{\alpha,\ell}$,
$c_{\beta,\ell}$, and $c_{\gamma,\ell}$ are initialized to zero,
so the initial activity probability equals $\rho$.
Earlier layers retain their previously selected parameters.

At each depth, we first train the new layer and adjust the variance
scaling factor used in the preceding layer's outgoing message. We then
fine-tune all layers jointly using a reduced learning rate. Both stages minimize
\eqref{eq:twc-layer-objective}. The checkpoint with the smallest
final-layer channel and activity loss on the validation set initializes
the next depth. Validation uses the training normalization constants.
After training to depth $L$, the selected layer parameters are fixed
for the receiver summarized in Section~\ref{sec:overall}.

\section{Overall Algorithm}
\label{sec:overall}

The updates in Section~\ref{sec:algorithm}, with parameters learned
as described in Section~\ref{sec:learning}, define the $L$-layer
MP receiver with a score-based channel prior. For $K$ candidate users, inference comprises $L$
linear updates and $LK$ evaluations of each network. Phase and noise
estimation first determine the observation model used by these layers.
We summarize this complete procedure and account for both its
parameter-estimation and unfolded-receiver costs.

\subsection{Message Scheduling}

From the received pilots, the receiver estimates the noise variance
by \eqref{eq:noise-mdl} and the phase increments by the fitting
procedure in Appendix~\ref{app:frontend}. It then constructs
$\widehat{\mathbf A}$ using \eqref{eq:effective-pilot}. The observations
and estimated noise variance are divided by $s$ and $s^2$,
respectively, as specified in Section~\ref{sec:system}. The resulting
$\mathbf A_{\rm c}$ and $\sigma_{\rm c}^2$ remain fixed throughout
the unfolded layers. The trained network parameters are also fixed;
the estimates from the frontend replace the sensing matrix and noise
variance supplied during training.

The first layer starts from the prior mean and variance
\begin{equation}
 \mathbf M_1=\mathbf0,\qquad w_1=w_{\rm ref}=\rho P_H.
 \label{eq:twc-initialization}
\end{equation}
Each layer updates the conditional channels and activity
probabilities, combines them into the effective channel estimate,
and returns an extrinsic message to the next layer. Let
$\widehat H_{k,L}$ denote the complex channel tensor obtained by
reshaping $\mu_{kL}$. After $L$ layers, we restore the original channel
scale and make the activity decisions as
\begin{equation}
 \begin{gathered}
 \widehat H_k^{\rm out}=s\widehat H_{k,L},\qquad
 \widehat X_k^{\rm out}=s p_{kL}\widehat H_{k,L},\\
 \widehat\alpha_k=\mathbf1\{p_{kL}\geq1/2\}.
 \end{gathered}
 \label{eq:twc-final-output}
\end{equation}
The activity probabilities $p_{kL}$ are also retained as soft outputs.
Table~\ref{tab:twc-algorithm} summarizes the complete receiver.

\begin{table}[t]
\centering
\caption{Deep-unfolded MP receiver with a score-based channel prior for JADCE}
\label{tab:twc-algorithm}
\begin{tabular}{p{0.93\linewidth}}
\toprule
\textbf{Input:} received pilots $\mathbf Y$, original pilot tensor $P$,
activity probability $\rho$, trained layer parameters, normalization
factor $s$, and average channel power $P_H$.\\
\textbf{Preprocessing:} estimate $\widehat\sigma^2$ and
$\widehat{\bm\beta}$ by Appendix~\ref{app:frontend}; construct
$\mathbf A_{\rm c}=\widehat{\mathbf A}$.
Normalize $\mathbf Y$ and $\widehat\sigma^2$ by $s$ and $s^2$,
and set $\sigma_{\rm c}^2$ to the normalized noise estimate.\\
\textbf{Initialize:} use \eqref{eq:twc-initialization} and compute
the thin SVD of $\mathbf A_{\rm c}$.\\
\textbf{For} $\ell=1,\ldots,L$:\\
1. Compute the linear update and the mean and variance passed to the denoiser using
\eqref{eq:twc-linear-correction}--\eqref{eq:twc-query-calibrated}.\\
2. Estimate the channels and activity probabilities, and combine
them using \eqref{eq:twc-conditional-mean}, \eqref{eq:twc-activity},
and \eqref{eq:twc-estimate}.\\
3. If $\ell<L$, form the feedback message and damp it using
\eqref{eq:twc-residual-scale}--\eqref{eq:twc-damping}.\\
\textbf{Return:} the channel estimates and decisions in
\eqref{eq:twc-final-output}, together with $p_{kL}$.\\
\bottomrule
\end{tabular}
\end{table}

\subsection{Computational Complexity Analysis}

The online cost consists of phase and noise estimation, followed by
the unfolded receiver. For the phase estimator in
Appendix~\ref{app:frontend}, let $G_{\rm c}$ be the number of coarse
delay--Doppler path matrices, $B_\phi\geq Q$ the number of phase bins,
and $S_\phi$ the number of coordinate sweeps. Precomputing the path
matrices and their pilot responses costs
$O(G_{\rm c}N^3+KG_{\rm c}QN^2)$.
Each user search requires correlations at cost $O(G_{\rm c}QNM)$
and phase-bin evaluation by FFTs at cost
$O(G_{\rm c}MB_\phi\log B_\phi)$.

Let $E_\phi$ count path-matrix evaluations when ranking Doppler candidates and during
local refinement, including derivative evaluations up to a constant
factor. Each evaluation, pilot response, and gain fit costs
$O(N^3+QN^2+QNM)$. For $K\leq QN$, a joint gain refit after each
sweep costs $O(QNK^2+QNKM)$. Including the initial user search,
the phase-estimation cost is
\begin{equation}
 \begin{aligned}
 C_{\rm phase}=O\big(&G_{\rm c}N^3+KG_{\rm c}QN^2\\
 &+(S_\phi+1)KG_{\rm c}M(QN+B_\phi\log B_\phi)\\
 &+E_\phi(N^3+QN^2+QNM)\\
 &+S_\phi(QNK^2+QNKM)\big).
 \end{aligned}
 \label{eq:twc-phase-complexity}
\end{equation}
Forming the spatial covariance and computing its eigenvalues for
noise estimation cost $O(QNM^2+M^3)$.

Within the unfolded receiver, the dominant computations are linear
estimation and network evaluation.
Since $\mathbf A_{\rm c}$ is fixed, one thin singular
value decomposition suffices to evaluate both the mean correction
and the trace in \eqref{eq:twc-linear-correction} across all layers.
With $r=\min(Q,J)$, this decomposition costs $O(QJr)$.
Each layer requires $O(QJC)$ operations for the residuals and
$O((Q+J)rC)$ for applying the singular vectors. If $F_s$ and $F_a$
denote the costs of one user's score-network and activity-classifier
evaluations, respectively, the cost per layer is
\begin{equation}
 C_{\rm layer}=O\left(QJC+(Q+J)rC+K(F_s+F_a)\right).
 \label{eq:twc-layer-complexity}
\end{equation}
The complete online cost is therefore
\begin{equation}
 O\left(C_{\rm phase}+QNM^2+M^3+QJr+L C_{\rm layer}\right).
 \label{eq:twc-complexity}
\end{equation}
For $Q\leq J$, the SVD and unfolded layers together cost
$O(Q^2J+L[QJC+K(F_s+F_a)])$, in addition to the phase and noise
estimators.
The depth $L$ therefore sets the number of repeated linear and neural
computations, while phase and noise estimation contribute a separate
cost. Residual-based variance estimation requires no second-order
score-network evaluation.
If the score network and activity classifier contain $P_s$ and $P_a$
parameters, the $L$ layers store $L(P_s+P_a+3)$ parameter scalars.
During inference, layers are evaluated sequentially. Excluding network
parameters and activations, the operator, SVD, and messages require
$O(QJ+(Q+J)r+JC+QC)$ storage. The phase estimator additionally
caches $O(G_{\rm c}N^2+KG_{\rm c}QN)$ entries, with
$O(G_{\rm c}B_\phi M)$ temporary FFT storage.

% \section{Simulation Results}
\label{sec:simulations}

% TODO: Write this section before enabling its input in main.tex.

% \section{Conclusion}
\label{sec:conclusion}

% TODO: Write this section before enabling its input in main.tex.


\appendices
\numberwithin{equation}{section}
\section{Phase and Noise Estimation}
\label{app:frontend}

\subsection{Inter-frame phase estimation}
Stack the received samples as
$[\mathbf Y_{\rm s}]_{qN+r,m}=Y_{qmr}$, so that
$\mathbf Y_{\rm s}\in\mathbb C^{QN\times M}$.
To estimate the phase increments, we approximate the contribution of each
candidate terminal by one dominant path with an unconstrained complex
gain at each antenna.
Absorbing the common scalar phase into the gain, we define the path matrix as
\begin{equation}
 \begin{aligned}
 \mathbf G_0(d,\nu)&=e^{-j2\pi c_1d^2}\mathbf G(d,\nu;0),\\
 [\mathbf h_k(d,\nu)]_{qN+r}
 &=e^{j2\pi\nu qT_{\rm frame}}
 \sum_{n=0}^{N-1}[\mathbf G_0(d,\nu)]_{rn}P_{qkn}.
 \end{aligned}
 \label{eq:phase-atom}
\end{equation}
The delay search is restricted to $d\in[0,N)$ and the Doppler search to $|\nu|\leq\nu_{\max}$. Under $2Nc_1\in\mathbb N_+$, shifting
the delay by an integer multiple of $N$ changes the path operator by
a scalar phase, which the fitted gains absorb. With
$\mathbf h_k=\mathbf h_k(d_k,\nu_k)$ and gain vector
$\mathbf b_k\in\mathbb C^M$, the fitting objective is
\begin{equation}
 \min_{\{d_k,\nu_k,\mathbf b_k\}}
 \left\|\mathbf Y_{\rm s}-\sum_{k=1}^K
 \mathbf h_k\mathbf b_k^{\mathsf T}\right\|_F^2.
 \label{eq:phase-fit}
\end{equation}
When updating terminal $k$, we fix the other contributions and define
$\mathbf Z_k=\mathbf Y_{\rm s}-\sum_{j\ne k}
\mathbf h_j\mathbf b_j^{\mathsf T}$. For a given $(d,\nu)$, the
least-squares gain and the resulting objective are given by
\begin{equation}
 \begin{aligned}
 \mathbf b_k^{\mathsf T}(d,\nu)
 &=\frac{\mathbf h_k(d,\nu)^{\mathsf H}\mathbf Z_k}
 {\|\mathbf h_k(d,\nu)\|^2},\\
 J_k(d,\nu)&=\|\mathbf Z_k\|_F^2-
 \frac{\|\mathbf h_k(d,\nu)^{\mathsf H}\mathbf Z_k\|^2}
 {\|\mathbf h_k(d,\nu)\|^2}.
 \end{aligned}
 \label{eq:phase-variable-projection}
\end{equation}
The gains are eliminated from the separable nonlinear least-squares
problem by variable projection \cite{oleary2013}. Coarse correlations over
delay, within-frame Doppler, and frame phase provide initial candidates.
A phase candidate $\beta$ corresponds to Doppler aliases
\begin{equation}
 \nu_j=\frac{\beta+2\pi j}{2\pi T_{\rm frame}},\qquad
 j\in\mathbb Z,\quad |\nu_j|\leq\nu_{\max}.
 \label{eq:phase-aliases}
\end{equation}
Candidates are ranked using the pilot response vector in
\eqref{eq:phase-atom}, which enforces the same Doppler within and across
frames. The selected candidates undergo bounded local refinement by
L-BFGS-B \cite{zhu1997}, with delay restricted to one differentiable
branch of the wrapped waveform. Each coordinate update is accepted
when it does not increase the residual energy, up to numerical tolerance.
After each sweep, all antenna gains are refitted jointly by least squares.

The receiver uses a fixed number of complete sweeps, retaining the previous
position as an additional refinement starting point after the first sweep.
It sets $\widehat\beta_k=\operatorname{wrap}
(2\pi\widehat\nu_kT_{\rm frame})$ and constructs
$\widehat{\mathbf A}$ by \eqref{eq:effective-pilot}.
This local fit provides phase estimates for all candidate terminals.
The nonlinear denoiser then estimates the full channel tensor in
\eqref{eq:channel}.

\subsection{Noise variance estimation}
Let $L=QN>M\geq2$ and form the spatial sample covariance
\begin{equation}
 \begin{aligned}
 \widehat{\mathbf R}_y&=\frac1L\sum_{q=0}^{Q-1}\sum_{r=0}^{N-1}
 \mathbf y_{qr}\mathbf y_{qr}^{\mathsf H},\\
 \mathbf y_{qr}&=(Y_{q0r},\ldots,Y_{q,M-1,r})^{\mathsf T}.
 \end{aligned}
 \label{eq:spatial-noise-covariance}
\end{equation}
Write its eigenvalues in descending order as
$\lambda_1\geq\cdots\geq\lambda_M$; values below
$\epsilon_{\rm mach}\lambda_1$ are floored at that level before
taking logarithms. For a candidate spatial signal rank $j$, define
\begin{equation}
 \bar\lambda_j=\frac1{M-j}\sum_{i=j+1}^M\lambda_i,\qquad
 \log g_j=\frac1{M-j}\sum_{i=j+1}^M\log\lambda_i.
 \label{eq:noise-tail-statistics}
\end{equation}
The Wax--Kailath minimum description length (MDL) criterion, as reviewed
in \cite{nadakuditi2008}, gives
\begin{equation}
 \begin{aligned}
 \widehat r_{\rm s}
 &=\arg\min_{0\leq j<M}\Big\{
 L(M-j)(\log\bar\lambda_j-\log g_j)\\
 &\hspace{6em}+\frac{j(2M-j)}2\log L\Big\},\\
 \widehat\sigma^2&=\bar\lambda_{\widehat r_{\rm s}}.
 \end{aligned}
 \label{eq:noise-mdl}
\end{equation}
Under a lower-rank spatial signal plus white-noise model, the criterion
selects the spatial signal rank and estimates the noise variance from
the remaining eigenvalues. Both $\widehat\sigma^2$ and
$\widehat{\mathbf A}$ remain fixed throughout the MP iterations.




\bibliographystyle{IEEEtran}
\bibliography{references}
\end{document}

